\scanpage{16}
\mspage{3}
revient au même, dans $W$ (puisque par hypothèse elle est déjà dans $C$). \deleted{\unsure{Or}}\unsure{Si} cette condition est satisfaite, alors $\varinjlim u_i:\varinjlim X_i\longrightarrow\varinjlim Y_i$ est dans $W$.

Dém. du corollaire. Supposons $u_.\in(TC)_I$, alors par la proposition, $\varinjlim_{i\in I}u_i$ est dans $TC$, \deleted{\unsure{donc dans $W$}} \inserted{ou ce qui revient} au $=$ dans $W$. D'autre part, appliquant ce $=$ résultat aux syst. inductifs induits sur $I_{<i_0}$, on trouve que dans $(\alpha)$, la \deleted{\unsure{$\beta$}} flèche \deleted{$\unsure{\lim}_{i<i_0}u_i$} $\alpha_{i_0}$ est dans $TC$, donc aussi $\beta_{i_0}$ par cochgt. de base, donc aussi $u_{i_0}=\widetilde{u}_{i_0}\beta_{i_0}$. \deleted{\unsure{Donc}} Donc il reste à prouver qu'inversement, si on suppose les $u_{i_0}\in W$, alors les $\widetilde{u}_{i_0}$ sont eux aussi dans $W$. (i.e. dans $TC=C\cap W$), ou ce qui revient au $=$, que les $\beta_{i_0}$ sont dans $W$, ou encore dans $TC$. Mais comme $TC$ est stable par coch. de base, il suffit \inserted{pour celà} de prouver que $\alpha_{i_0}\in TC$, ou encore $\alpha_{i_0}\in C$. Donc on est ramené à prouver que \deleted{\unsure{si} $u_.\in W_{I^\circ}$}
\[
 u_.\in W_{I^\circ}\cap C_I\Longrightarrow \left(\varinjlim_{\unsure{i}<i_0}u_i\right)\in W \tag*{$c_{i_0})$}
\]
pour \unsure{tt} $i_0\in I$, \unsure{ou}, changeant de notations, à prouver

Nous prouverons ensemble a) b) c) par induction transfinie sur le nb ordinal $\alpha$, tel que $I\simeq I_\alpha=\{\text{ens des ordinaux }i<\alpha\}$.
% archive sheet number: 15

